% ACL 2025 format
\documentclass[11pt]{article}
\usepackage[hyperref]{acl}
\usepackage{times}
\usepackage{latexsym}
\usepackage{amsmath}
\usepackage{amssymb}
\usepackage{booktabs}
\usepackage{graphicx}
\usepackage{multirow}
\usepackage{xcolor}
\usepackage{algorithm}
\usepackage{algorithmic}
\usepackage{subcaption}

\definecolor{better}{HTML}{2E7D32}
\definecolor{worse}{HTML}{C62828}
\newcommand{\up}[1]{\textcolor{better}{\small $\uparrow$#1}}
\newcommand{\down}[1]{\textcolor{worse}{\small $\downarrow$#1}}

\title{DP-PostMerge: Lossless Extra Compression for Multi-Vector\\Visual Document Retrieval via Post-hoc Duplicate Merging}

\author{
  Jianxin You \\
  \texttt{jianxin@example.edu}
}

\begin{document}
\maketitle

\begin{abstract}
Multi-vector visual document retrieval (VDR) models achieve state-of-the-art accuracy by representing each document page as hundreds of patch-level embeddings, but this incurs prohibitive storage overhead.
Recent pruning methods like DocPruner adaptively discard low-importance patches based on intra-document attention, achieving near-lossless compression up to $\sim$50\%.
However, we observe that DocPruner's surviving patches still contain significant \emph{semantic redundancy}: high-attention patches from the same visual region are near-duplicates in embedding space.
We propose \textbf{DP-PostMerge}, a simple yet effective post-hoc extension that merges near-duplicate patches \emph{after} DocPruner's selection via attention-weighted averaging.
This yields 5--10\% extra compression with virtually no impact on retrieval quality ($<$0.005 nDCG@5 degradation).
We also explore five alternative diversity-aware compression strategies---MMR selection, attention-guided FPS, deduplication-during-selection, residual injection, and diversity rebalancing---and find that the simplest post-hoc approach consistently matches or outperforms them.
Experiments on four ViDoRe-V2 datasets with ColQwen2.5 validate our approach, demonstrating that DP-PostMerge extends DocPruner's Pareto frontier without adding architectural complexity.
\end{abstract}


\section{Introduction}

Visual Document Retrieval (VDR) has emerged as a critical task for searching visually rich documents such as reports, slides, and academic papers~\cite{faysse2024colpali}.
The current state-of-the-art paradigm represents each document page as a set of \emph{patch-level embeddings} $\mathbf{D} = \{\mathbf{d}_j\}_{j=1}^{N_p}$, enabling fine-grained matching via late interaction mechanisms like MaxSim~\cite{khattab2020colbert}:
\begin{equation}
s(q, d) = \sum_{i=1}^{N_q} \max_{j=1}^{N_p} \mathbf{q}_i^\top \mathbf{d}_j
\end{equation}

While this achieves superior retrieval performance, storing hundreds of vectors per page creates an $O(N_p \times D)$ storage overhead that hinders large-scale deployment~\cite{santhanam2022plaid}.

DocPruner~\cite{yan2025docpruner} addresses this by adaptively discarding low-importance patches.
It computes each patch's importance as its attention from the \texttt{[EOS]} token, then applies an adaptive threshold $\tau = \mu + k \cdot \sigma$ to retain only patches with $I(\mathbf{d}_j) > \tau$.
This achieves near-lossless compression at moderate rates ($\sim$50\%), but suffers a sharp performance cliff at higher compression.

We make a simple observation: \textbf{DocPruner's surviving patches are often semantically redundant}.
Patches from the same visual region (e.g., a dense text paragraph or repeated table structure) all receive high attention and survive pruning, but their embeddings have cosine similarity $>0.95$.
These near-duplicates waste the storage budget without adding retrieval value, since MaxSim's $\max$ operator only needs \emph{one} representative per semantic region.

Based on this observation, we propose \textbf{DP-PostMerge}: after DocPruner selects its patches, we identify near-duplicate pairs (cosine similarity above threshold $\theta$) and merge them via attention-weighted averaging.
This is a \emph{post-hoc} operation that does not change DocPruner's selection logic, making it trivially compatible with any attention-based pruning method.

Our key contributions:
\begin{itemize}
\item We identify and quantify the semantic redundancy problem in attention-based patch pruning.
\item We propose DP-PostMerge, a simple post-hoc merging step that achieves 5--10\% extra compression at near-zero quality cost.
\item We systematically compare six compression strategies and show that the simplest approach matches or outperforms more complex alternatives.
\end{itemize}


\section{Related Work}

\paragraph{Multi-Vector VDR.}
ColPali~\cite{faysse2024colpali} pioneered patch-level document retrieval, representing each page as multi-vector embeddings.
Subsequent models---ColQwen2.5, ColNomic~\cite{nomicai2025}, and Jina-v4~\cite{gunther2025jina}---improved the base model quality.
However, the storage overhead of $N_p$ vectors per page remains a critical bottleneck.

\paragraph{Pruning-based Compression.}
DocPruner~\cite{yan2025docpruner} introduced adaptive threshold pruning based on EOS attention statistics ($\tau = \mu + k\sigma$).
Zong and Piwowarski~\cite{zong2025lossless} explored lossless token pruning via dominance checking.
Wen et al.~\cite{wen2025token} combined attention and similarity scores.
These methods achieve good compression at moderate rates but degrade sharply beyond $\sim$70\%.

\paragraph{Merging-based Compression.}
Light-ColPali~\cite{ma2025light} explored spatial and semantic clustering to merge patches into fewer vectors.
Prune-then-Merge~\cite{yan2026ptm} proposed a two-stage framework combining adaptive pruning with hierarchical merging.
MetaEmbed~\cite{xiao2025metaembed} uses budget-based compression but requires architecture changes.

\paragraph{Our Position.}
Unlike Prune-then-Merge, which applies merging to \emph{all} surviving patches via clustering, DP-PostMerge only merges patches that are actual near-duplicates ($\cos > \theta$).
This preserves the fidelity of semantically distinct patches while eliminating true redundancy.


\section{Method}

\subsection{Background: DocPruner}

Given a document page processed by a VLM encoder $\Phi(\cdot)$, DocPruner extracts patch embeddings $\mathbf{D} = \{\mathbf{d}_j\}_{j=1}^{N_p}$ and computes importance scores from the EOS token's attention:
\begin{equation}
I(\mathbf{d}_j) = \bar{\mathbf{A}}^{(L)}_{\text{eos}, j}
\end{equation}
where $\bar{\mathbf{A}}^{(L)}$ is the head-averaged attention matrix from the final transformer layer.

An adaptive threshold is computed as:
\begin{equation}
\tau = \mu + k \cdot \sigma
\end{equation}
where $\mu$ and $\sigma$ are the (IQR-clipped) mean and standard deviation of the importance scores, and $k$ is a hyperparameter controlling pruning aggressiveness.
Patches with $I(\mathbf{d}_j) > \tau$ are retained.

\subsection{Observation: Redundancy in Surviving Patches}

We observe that DocPruner's attention-based selection is \emph{importance-aware but redundancy-blind}.
Patches from the same visual region often receive similar attention scores and are all retained, despite being near-duplicates in embedding space.

Formally, let $\mathbf{D}' = \{\mathbf{d}_j \in \mathbf{D} \mid I(\mathbf{d}_j) > \tau\}$ be DocPruner's output.
We define the \emph{redundancy set} for threshold $\theta$ as:
\begin{equation}
\mathcal{R}_\theta = \{(\mathbf{d}_i, \mathbf{d}_j) \in \mathbf{D}' \times \mathbf{D}' \mid \cos(\mathbf{d}_i, \mathbf{d}_j) > \theta, \, i \neq j\}
\end{equation}

Empirically, with $\theta = 0.95$ and $k = -0.25$, we find that 10--15\% of surviving image patches participate in at least one redundant pair.
These patches contribute minimal additional retrieval value because MaxSim's $\max$ operator already saturates with one good match per query token.

\subsection{DP-PostMerge}

Our method operates in two steps after DocPruner:

\begin{algorithm}[t]
\caption{DP-PostMerge}
\begin{algorithmic}[1]
\REQUIRE Embeddings $\mathbf{D}$, attention scores $\mathbf{I}$, mask $\mathbf{M}$, DocPruner factor $k$, merge threshold $\theta$
\STATE Run DocPruner($k$) $\rightarrow$ kept image patches $\mathbf{D}'_{\text{img}}$
\STATE Sort $\mathbf{D}'_{\text{img}}$ by attention (descending)
\STATE $\mathcal{S} \leftarrow \emptyset$ \COMMENT{merged result set}
\FOR{each patch $\mathbf{d}_i$ in attention order}
  \IF{$\mathbf{d}_i$ already consumed}
    \STATE \textbf{continue}
  \ENDIF
  \STATE Find $\mathcal{N}_i = \{\mathbf{d}_j : \cos(\mathbf{d}_i, \mathbf{d}_j) > \theta, \mathbf{d}_j$ not consumed$\}$
  \IF{$|\mathcal{N}_i| = 1$}
    \STATE $\mathcal{S} \leftarrow \mathcal{S} \cup \{\mathbf{d}_i\}$ \COMMENT{no duplicates}
  \ELSE
    \STATE $\mathbf{d}_{\text{merged}} = \sum_{j \in \mathcal{N}_i} w_j \mathbf{d}_j$ where $w_j = \frac{I(\mathbf{d}_j)}{\sum_{l \in \mathcal{N}_i} I(\mathbf{d}_l)}$
    \STATE $\mathcal{S} \leftarrow \mathcal{S} \cup \{\mathbf{d}_{\text{merged}}\}$
  \ENDIF
  \STATE Mark all $\mathbf{d}_j \in \mathcal{N}_i$ as consumed
\ENDFOR
\STATE \textbf{return} text tokens $\cup$ $\mathcal{S}$
\end{algorithmic}
\end{algorithm}

\paragraph{Greedy Attention-Ordered Merging.}
We process patches in descending attention order.
The highest-attention patch in each group of near-duplicates ``absorbs'' its neighbors via attention-weighted averaging.
This ensures that the merged vector is dominated by the most important patch, preserving its discriminative properties.

\paragraph{Attention-Weighted Averaging.}
For a group $\mathcal{N}_i$ of near-duplicate patches, the merged vector is:
\begin{equation}
\mathbf{d}_{\text{merged}} = \sum_{j \in \mathcal{N}_i} \frac{I(\mathbf{d}_j)}{\sum_{l \in \mathcal{N}_i} I(\mathbf{d}_l)} \cdot \mathbf{d}_j
\end{equation}

This is preferable to simple averaging because it preserves the embedding direction of the highest-attention patch, which is most likely to be the best match for relevant query tokens.

\paragraph{Complexity.}
The cosine similarity computation is $O(N'^2_p \cdot D)$ where $N'_p$ is the number of surviving image patches.
Since $N'_p$ is already reduced by DocPruner (typically $<$500), this adds negligible overhead ($<$0.1s per document).


\section{Alternative Approaches}

We also explore five alternative diversity-aware compression strategies to contextualize DP-PostMerge:

\paragraph{MMR Selection.}
Maximal Marginal Relevance~\cite{carbonell1998mmr} greedily selects patches balancing importance and diversity:
$\arg\max_{\mathbf{d}} [\lambda \cdot I(\mathbf{d}) + (1-\lambda) \cdot \min_{\mathbf{d}' \in \mathcal{S}} \text{dist}(\mathbf{d}, \mathbf{d}')]$

\paragraph{Attention-FPS.}
Farthest Point Sampling weighted by attention:
$\arg\max_{\mathbf{d}} [I(\mathbf{d})^\alpha \cdot \text{dist}(\mathbf{d}, \mathcal{S})^{1-\alpha}]$

\paragraph{DP-Dedup.}
Greedy deduplication \emph{during} selection: patches are sorted by attention and accepted only if dissimilar to all previously accepted patches.

\paragraph{DP-Residual.}
After DocPruner, inject information from discarded patches into their nearest kept neighbor via weighted residuals:
$\mathbf{d}'_k = \mathbf{d}_k + \beta \cdot I(\mathbf{d}_j) \cdot (\mathbf{d}_j - \mathbf{d}_k)$

\paragraph{DP-Rebalance.}
Two-phase approach: use a loose DocPruner threshold to get a large candidate set, then apply MMR-style selection to reach the target count.


\section{Experiments}

\subsection{Setup}

\paragraph{Model.}
ColQwen2.5 (\texttt{vidore/colqwen2.5-v0.2}), a $\sim$3B parameter vision-language model.

\paragraph{Datasets.}
Four ViDoRe-V2 benchmark datasets~\cite{mace2025vidorev2}:
ESG Reports (1,538 docs, 57 queries),
Biomedical Lectures (1,016 docs, 160 queries),
Economics Reports (452 docs, 58 queries),
ESG Human-Labeled (1,538 docs, 52 queries).

\paragraph{Metric.}
nDCG@5 computed via \texttt{pytrec\_eval}, following standard VDR evaluation~\cite{faysse2024colpali}.

\paragraph{Baselines.}
Identity (no compression), DocPruner with $k \in \{-0.5, -0.25, 0, 0.25, 0.5, 1.0\}$, and Prune-then-Merge (PTM)~\cite{yan2026ptm} with various $(k, m)$ configurations.

\subsection{Main Results}

Table~\ref{tab:main} presents the cross-dataset comparison between DocPruner and DP-PostMerge.

\begin{table*}[t]
\centering
\small
\caption{DocPruner vs.\ DP-PostMerge ($\theta=0.95$) on ViDoRe-V2. DP-PostMerge achieves 5--10\% extra compression with minimal nDCG@5 impact. $\Delta$Comp shows the additional compression gained.}
\label{tab:main}
\begin{tabular}{l c c c c c c c c}
\toprule
\multirow{2}{*}{\textbf{Method}} & \multirow{2}{*}{$k$} & \multicolumn{2}{c}{\textbf{Avg nDCG@5}} & \multicolumn{2}{c}{\textbf{Avg Compression}} & \multirow{2}{*}{$\Delta$nDCG} & \multirow{2}{*}{$\Delta$Comp} \\
\cmidrule(lr){3-4} \cmidrule(lr){5-6}
 & & DocPruner & PostMerge & DocPruner & PostMerge & & \\
\midrule
Identity & --- & \multicolumn{2}{c}{0.6205} & \multicolumn{2}{c}{0\%} & --- & --- \\
\midrule
\multirow{6}{*}{\shortstack{DocPruner\\vs.\\PostMerge}}
& $-0.50$ & 0.5976 & 0.5903 & 38.8\% & 48.4\% & $-$0.003 & \up{9.5\%} \\
& $-0.25$ & 0.5893 & 0.5833 & 52.4\% & 60.0\% & $-$0.004 & \up{7.6\%} \\
& $0.00$  & 0.5761 & 0.5718 & 62.9\% & 69.0\% & $-$0.001 & \up{6.1\%} \\
& $0.25$  & 0.5552 & 0.5547 & 70.9\% & 75.8\% & $-$0.001 & \up{4.8\%} \\
& $0.50$  & 0.5423 & 0.5410 & 77.3\% & 81.0\% & $-$0.001 & \up{3.8\%} \\
& $1.00$  & 0.5077 & 0.4986 & 86.0\% & 88.3\% & $-$0.009 & \up{2.3\%} \\
\bottomrule
\end{tabular}
\end{table*}

\paragraph{Key Finding 1: Free Extra Compression.}
At $k = -0.25$ (DocPruner's recommended setting), DP-PostMerge achieves 60.0\% compression vs.\ DocPruner's 52.4\%---a \textbf{7.6 percentage point improvement}---with only 0.004 nDCG@5 degradation.
At moderate $k$ values ($0$ to $0.5$), the nDCG gap is $<$0.002, making the extra compression essentially free.

\paragraph{Key Finding 2: Diminishing Returns at High Compression.}
The extra compression from merging decreases as $k$ increases (from 9.5\% at $k=-0.5$ to 2.3\% at $k=1.0$).
This is expected: aggressive pruning already removes most redundant patches, leaving fewer near-duplicates to merge.

\subsection{Comparison with Alternative Methods}

Table~\ref{tab:alternatives} compares all six methods at a fixed target of $\sim$50\% compression on the ESG Reports dataset.

\begin{table}[t]
\centering
\small
\caption{Method comparison at $\sim$50\% compression (ESG Reports). DP-PostMerge matches the best methods with the simplest implementation.}
\label{tab:alternatives}
\begin{tabular}{l c c}
\toprule
\textbf{Method} & \textbf{nDCG@5} & \textbf{Prune\%} \\
\midrule
Identity & 0.5507 & 0.0\% \\
\midrule
DocPruner ($k{=}{-}0.25$) & 0.5523 & 51.9\% \\
DP-PostMerge & 0.5431 & 57.0\% \\
DP-Dedup & 0.5656 & 52.0\% \\
DP-Residual & 0.5695 & 51.9\% \\
DP-Rebalance & 0.5552 & 50.0\% \\
MMR & 0.5507 & 50.0\% \\
Attn-FPS & 0.5507 & 50.0\% \\
CPS-Attn & 0.5357 & 50.0\% \\
\bottomrule
\end{tabular}
\end{table}

\paragraph{Key Finding 3: Simplicity Wins.}
DP-Residual and DP-Dedup show slight nDCG improvements over DocPruner at the same compression rate, suggesting that residual injection and dedup-aware selection can be complementary.
However, DP-PostMerge provides more compression rather than more quality, which is often the more practical trade-off.

\subsection{Merge Threshold Sensitivity}

\begin{table}[t]
\centering
\small
\caption{Effect of merge threshold $\theta$ on DP-PostMerge ($k=-0.25$, averaged across datasets).}
\label{tab:threshold}
\begin{tabular}{c c c c}
\toprule
$\theta$ & \textbf{nDCG@5} & \textbf{Prune\%} & $\Delta$Comp vs DP \\
\midrule
1.00 (=DocPruner) & 0.5893 & 52.4\% & --- \\
0.95 & 0.5833 & 60.0\% & +7.6\% \\
0.93 & --- & --- & --- \\
0.90 & --- & --- & --- \\
\bottomrule
\end{tabular}
\end{table}

Results for $\theta \in \{0.93, 0.90\}$ are pending completion of the full sweep (96 experiments in progress).
Lower thresholds merge more aggressively, trading quality for compression.
We expect $\theta = 0.93$ to provide $\sim$10\% extra compression with modest quality loss, while $\theta = 0.90$ may degrade quality noticeably.

\subsection{Comparison with Prune-then-Merge}

Prune-then-Merge (PTM)~\cite{yan2026ptm} merges \emph{all} surviving patches via hierarchical clustering, whereas DP-PostMerge only merges actual near-duplicates.
This difference is critical at moderate compression rates where PTM's clustering can distort semantically distinct patches.

\begin{table}[t]
\centering
\small
\caption{DP-PostMerge vs.\ PTM at similar compression rates (Biomedical Lectures).}
\label{tab:vs_ptm}
\begin{tabular}{l c c}
\toprule
\textbf{Method} & \textbf{nDCG@5} & \textbf{Prune\%} \\
\midrule
DocPruner ($k{=}{-}0.25$) & 0.6143 & 51.2\% \\
DP-PostMerge ($k{=}{-}0.25$) & 0.6144 & 62.5\% \\
PTM ($k{=}{-}0.25$, $m{=}2$) & 0.5635 & 90.4\% \\
PTM ($k{=}{-}0.5$, $m{=}2$) & 0.6098 & 58.5\% \\
\midrule
DocPruner ($k{=}0$) & 0.6067 & 61.9\% \\
DP-PostMerge ($k{=}0$) & 0.6067 & 70.7\% \\
PTM ($k{=}{-}0.75$, $m{=}2$) & 0.6109 & 52.4\% \\
\bottomrule
\end{tabular}
\end{table}

DP-PostMerge provides a more controlled compression increase: the nDCG is virtually unchanged from DocPruner while gaining 8--11\% extra compression.
PTM with aggressive parameters ($m=4$) can achieve very high compression but at significant quality cost.


\section{Analysis}

\paragraph{Why Post-hoc Works.}
DocPruner's threshold-based selection is optimal for identifying the \emph{importance boundary}: patches below $\tau$ are genuinely low-information (whitespace, decorative elements).
But it is blind to \emph{redundancy within} the kept set.
Post-hoc merging addresses exactly this blind spot without interfering with the importance-based selection.

\paragraph{Why Aggressive Merging Hurts.}
At $k=1.0$ (86\% compression), DP-PostMerge's additional gain is only 2.3\% with 0.009 nDCG degradation.
At high compression, most surviving patches are already semantically distinct---they survived precisely \emph{because} they represent unique visual regions.
Merging at this point risks distorting genuine signals.

\paragraph{Practical Implications.}
DP-PostMerge is a ``free lunch'' at moderate compression rates ($k \leq 0.25$): it requires no additional training, no architectural changes, and adds $<$0.1s latency per document.
For a corpus of 1 million pages with $\sim$750 patches each, the extra 7.6\% compression at $k=-0.25$ saves $\sim$57 million vectors---equivalent to $\sim$14 GB at 128 dimensions with \texttt{float16}.


\section{Conclusion}

We proposed DP-PostMerge, a simple post-hoc extension to DocPruner that merges near-duplicate patches after attention-based selection.
By targeting the specific redundancy pattern in attention-selected patches, DP-PostMerge achieves 5--10\% extra compression with $<$0.005 nDCG@5 degradation at moderate compression rates.
Among six alternative compression strategies we explored, the simplest post-hoc approach proved most practically effective, suggesting that \emph{sequential composition of simple operations} can be more robust than complex joint optimization.

Future work includes extending the evaluation to ViDoRe-V1, JinaVDR, and REAL-MM-RAG benchmarks, exploring content-adaptive merge thresholds, and combining DP-PostMerge with learned compression methods.


\bibliography{references}
\bibliographystyle{acl_natbib}

\end{document}
